
\setcounter{page}{160}
\section{The duality theorem for projective space}

The duality theorem for projective space now follows easily from what has gone before. At the same time it is a model of how the duality is defined in terms of the functor \(f^!\) and the isomorphism \(\operatorname{Tr}_f\). When we have a satisfactory functorial theory of \(f^!\) (which is \(f^*\otimes\omega[n]\) in the smooth case) and \(\operatorname{Tr}_f\), we will prove the most general duality theorem by reducing to this case (Chapter VII).

Let \(Y\) be a noetherian prescheme of finite Krull dimension, let \(X=\mathbb{P}_{Y}^{n}\), and let \(f\colon X \to Y\) be the projection. We define the duality morphism
\[
\underline{\theta}_{f}\colon 
\Rfst \RHom_X\bigl(F^\bullet,f^! G^\bullet\bigr)
\to \RHom_Y\bigl(\Rfst F^\bullet,G^\bullet\bigr)
\]
for \(F^\bullet \in D^-(X)\) and \(G^\bullet \in D_{\mathrm{qc}}^+(Y)\) by composing the morphism [II.5.5] with the trace morphism in the second variable (Proposition 4.3).

Applying the functor \(\mathbf{R}\Gamma(Y,\,\cdot)\) to both sides, and using [II.5.2] and [II.5.3] we obtain a global duality morphism
\[
\theta_{f}\colon 
\Hom_{D(X)}\bigl(F^\bullet,f^! G^\bullet\bigr)
\to \Hom_{D(Y)}\bigl(\Rfst F^\bullet,G^\bullet\bigr),
\]

\setcounter{page}{161}
and taking the cohomology of this, we get morphisms
\[
\theta_{f}^{i}\colon 
\operatorname{Ext}_{X}^{i}\bigl(F^\bullet,f^! G^\bullet\bigr)
\to \operatorname{Ext}_{Y}^{i}\bigl(\Rfst F^\bullet,G^\bullet\bigr).
\]

\begin{theorem}
Let \(Y\) be a noetherian prescheme of finite Krull dimension, \(X=\mathbb{P}_{Y}^{n}\), and \(f\colon X\to Y\) the projection. Then the duality morphisms \(\underline{\theta}_f,\theta_f,\theta_f^i\) are isomorphisms for all \(F^\bullet \in D_{\mathrm{qc}}^-(X)\) and \(G^\bullet \in D_{\mathrm{qc}}^+(Y)\).
\end{theorem}

\begin{proof}
The question is local on \(Y\), so we may assume \(Y\) affine. Then every quasi-coherent sheaf on \(Y\) is a quotient of a direct sum of copies of \(\sheaf{O}_Y\); hence, as in Lemma 4.1, every quasi-coherent \(\sheaf{O}_X\)-module is a quotient of a direct sum of copies of \(\sheaf{O}_X(-m)\) for various \(m\). Using the isomorphism \(\omega \cong \sheaf{O}_X(-n-1)\), we see that any quasi-coherent sheaf \(\sheaf{F}\) on \(X\) is a quotient of a sheaf of the form
\[
L=\bigoplus \omega(-m_i)
\]
for integers \(m_i>0\).

The functors in question are way-out right in both variables, so by the Lemma on Way-out Functors [I.7.1], we reduce to the case \(F^\bullet = L\) and \(G^\bullet\) a single injective quasi-coherent sheaf.

\setcounter{page}{162}
Furthermore, \(\RHom\) transforms direct sums in the first variable to direct products, so we reduce to \(F^\bullet=\omega(-m)\) for \(m>0\). We must prove the map
\[
\begin{aligned}
\underline{\theta}_{f}\colon
&\Rfst \RHom_X\bigl(\omega(-m),f^!(G)\bigr) \\
&\to \RHom_Y\bigl(\Rfst(\omega(-m)),G\bigr)
\end{aligned}
\]
is an isomorphism.

By [II.5.16] and the projection formula, the left-hand side becomes
\[
\Rfst\bigl(\sheaf{O}_X(m)\bigr) \otimes G[n],
\]
since \(m>0\) and \(R^i f_*(\sheaf{O}_X(m))=0\) for \(i>0\) (Theorem 3.4). On the other hand, \(R^i f_*(\omega(-m))=0\) for \(i\neq n\), so the right-hand side becomes
\[
\Hom_Y\bigl(R^n f_*(\omega(-m)),G\bigr)[n],
\]
using again [II.5.16] and the fact that \(R^n f_*(\omega(-m))\) is locally free of finite rank on \(Y\).

Now \(\underline{\theta}_f\) is induced by the cup-product pairing
\[
f_*\bigl(\sheaf{O}_X(m)\bigr) \times R^n f_*\bigl(\omega(-m)\bigr)
\to R^n f_*(\omega),
\]
which is an isomorphism by Theorem 3.4.
\q.e.d.
\end{proof}

\setcounter{page}{163}
\begin{corollary}
Let \(A\) be a noetherian ring, \(X=\mathbb{P}_{A}^{n}\), let \(F^\bullet\in D_{\mathrm{qc}}^-(X)\), and let \(I\) be an injective \(A\)-module. Then there is a canonical isomorphism
\[
\Hom_A\bigl(H^{i}(X, F^\bullet), I\bigr)
\cong \operatorname{Ext}_{\sheaf{X}}^{n-i}\bigl(F^\bullet, \omega \otimes_A I\bigr).
\end{corollary}

\begin{remark}
When \(A\) is a field and \(I=A\), and \(F^\bullet\) is a single sheaf concentrated in degree zero, one recovers the classical Serre duality theorem for projective space over a field.
\end{remark}